\section{生产系统中的脆弱点：故障风暴与恢复设计}

官方 deck 给出平台目标，课堂口头案例则展示这些目标为何必要。本节从一次大型 chatbot 恢复风暴出发，分析冷启动、队列、流量放大和副本预热之间的反馈回路，并将其映射到多步 Agent 任务的 admission control 与降级设计。

\subsection{大型 Chatbot 事故案例}
故障发生后最直觉的动作是重启服务，但如果积压请求立即压向少量冷副本，恢复动作本身会制造下一次故障。本小节先重建这一因果链，再说明为什么恢复阶段必须限制放量速度，而不是只盯着副本是否重新上线。
在约 48 分钟，讲者分享了一个真实事故：客户误关闭主模型服务后，系统重启触发请求堆积，首个恢复副本被突发流量打死，随后第二个副本重复失败，形成典型恢复风暴（restart storm）。

\begin{importantbox}{恢复阶段往往比故障发生阶段更危险}
系统失效时损失可见；系统恢复时如果没有流控与预热策略，容易出现 ``刚恢复就再次崩溃'' 的连锁反应，导致停机时间显著拉长。
\end{importantbox}

\subsection{15 分钟恢复的关键动作}
该案例中团队通过在服务前加流量调节（traffic regulator），限制早期放量速度，让副本逐步预热，最终在约 15 分钟内恢复全链路。这是典型的 ``以吞吐爬坡换稳定恢复''。

\begin{knowledgebox}{可迁移的 SRE 经验}
\begin{itemize}
    \item 冷启动阶段必须有 admission control；
    \item 队列长度与并发阈值要联动，不可静态配置；
    \item 恢复流程要脚本化，避免临场手工操作放大错误；
    \item 关键组件需具备 brownout 模式，先保核心能力再恢复完整能力。
\end{itemize}
\end{knowledgebox}

\begin{warningbox}{Agent 系统的放大效应}
Agent 的请求通常是多步链路，一个环节失控会放大到整条工作流。传统 API 服务可承受的短抖动，在 Agent 系统里可能演化为任务雪崩。
\end{warningbox}

\subsection{本章小结}
生产系统的核心能力不是 ``平时跑得快''，而是 ``故障时不扩散、恢复时不复发''。Agent 场景下这一要求更高。

